\SourceLeaf{423}%
by the $3$ points $(1),(2),(3)$, which are therefore pairwise $90^{0}$ apart and are also to indicate the sense in which the coordinates are regarded as increasing. We collect here some familiar propositions that will be used throughout.

1) The angle between two intersecting straight lines is measured by the arc [of the great circle] between the points on the sphere that represent their directions.

2) The position of any plane can be represented on the sphere by the great circle in which the sphere is cut by a plane drawn through its center parallel to the original plane.

3) The angle between two planes is equal to the angle between the great circles that represent their positions and is therefore also measured by the angle between the poles of the great circles.

4) If $x,y,z$ and $x',y',z'$ are the coordinates of two points, $r$ their distance apart, and $L$ the point on the sphere that represents the direction of the straight line from the first point to the second, then
\[
\begin{aligned}
x'&=x+r\cos(1)L\\
y'&=y+r\cos(2)L\\
z'&=z+r\cos(3)L.
\end{aligned}
\]
\par\noindent\hspace*{1.2em}5) It follows readily that we always have
\[
\cos(1)L^{2}+\cos(2)L^{2}+\cos(3)L^{2}=1
\]
[and] also, if $L'$ is any other point on the sphere,
\[
\cos(1)L.\cos(1)L'+\cos(2)L.\cos(2)L'+\cos(3)L.\cos(3)L'=\cos LL'.
\]

We add another theorem that, to our knowledge, occurs nowhere else and can often be used to advantage.

Let $L,L',L'',L'''$ be four points on the sphere and $A$ the angle formed by $LL'''$ and $L'L''$ at their point of intersection. [Then we have]
\[
\cos LL'.\cos L''L'''-\cos LL''.\cos L'L'''=\sin LL'''.\sin L'L''.\cos A.
\]
